\clearpage
\fancyhead[L]{\small CONNAISSANCE DES TEMS.}
\SourceLeaf{620}
\begin{center}\small Göttingische gelehrte Anzeigen. Stück 191. Seite 1905..1910. 1819 November 29.\\[3pt]\rule{.72\linewidth}{.3pt}\end{center}

\noindent\emph{Connaissance des tems, ou des mouvemens célestes, à l'usage des astronomes et des navigateurs, pour l'an {\upshape 1821}. Publiée par le bureau des longitudes.} Paris 1818. Bei der Witwe Courcier. 336 S. gr. Octav.

Der Anhang des vorliegenden Jahrganges der Connaissance des tems zeichnet sich durch mehrere gehaltreiche Aufsätze von Frankreichs ersten Geometern vorzüglich aus. Zuerst Resultate einer theoretischen Untersuchung über die Libration des Mondes von \textsc{Poisson}. Es werden hier auch neuere Beobachtungen des Mondsflecken Manilius, von \textsc{Bouvard}, erwähnt, aus denen die Neigung des Mondsäquators $1^\circ\,27'\,40''$ folge (\textsc{Tobias Mayer} hatte 1749 aus Beobachtungen desselben Flecken $1^\circ\,29'$ gefunden). Die Vervielfältigung ähnlicher Beobachtungen wäre unstreitig sehr zu wünschen. --- Neue Methode, die Cometenbahnen zu bestimmen von \textsc{Lagrange}. Diese Abhandlung war schon in dem Berliner astronomischen Jahrbuch für 1783 deutsch erschienen. Die Methode setzt sechs vollständige Beobachtungen voraus, die paarweise nahe bei einander liegen müssen. Da, wenn die Lage der Ebene der Bahn bekannt ist, aus dem geocentrischen Orte der heliocentrische sich ableiten lässt, und zwei heliocentrische, einander nahe liegende Oerter bekanntlich eine sehr einfache Bestimmung des halben Parameters geben, so erhellet, dass jedes Paar einander nahe liegender geocentrischer Oerter eine Bedingungsgleichung zwischen der Neigung der Bahn, der Länge des aufsteigenden Knotens und dem halben Parameter gibt; die sechs vollständigen Beobachtungen geben daher drei solcher Bedingungsgleichungen, und durch die Elimination von zweien der unbekannten Grössen gelangt man denn endlich zu einer Gleichung des siebenten Grades. Bei dem heutigen Zustande dieses Theils des astronomischen Calculs wird man keine Veranlassung haben, von dieser Methode Gebrauch zu machen. --- Ueber die Rotation der Erde, von \textsc{Laplace}. Dieser Aufsatz enthält mehrere interessante Entwicklungen über die freie Drehungsaxe des Erdkörpers, in so fern er mit einer an ihrer Oberfläche im Gleichgewicht befindlichen Flüssigkeit bedeckt ist, und über die Veränderung der Lage von jener durch die Anziehung des Mondes und der Sonne. --- Den letzten Gegenstand behandelt gleichfalls auf ähnliche Weise die hierauf folgende Abhandlung von \textsc{Poisson} über die Präcession der Nachtgleichen mit der Concinnität, die man an diesem trefflichen Geometer gewohnt ist. --- Ueber den Einfluss der grossen Gleichung des Jupiter und Saturn auf die Bewegung der Körper des Sonnensystems von \textsc{Laplace}. Nur bei den eignen Trabanten jener Planeten wird dieser Einfluss merklich gefunden, wo er sich in der Gestalt einer Ungleichheit von ähnlicher Form, wie die grossen Gleichungen selbst offenbart: beim vierten Jupiterstrabanten findet \textsc{Laplace} den grössten Werth 14 Secunden. --- Geographische Lage einer grossen Anzahl von Punkten des mittelländischen Meeres, die in den Jahren 1816 und 1817 von dem Fregattencapitain \textsc{Gauttier} bestimmt sind. --- Ueber das Gesetz der Schwere, wenn das Erdsphäroid als homogen und von gleicher Dichtigkeit wie das Meer angenommen wird, von \textsc{Laplace}. Es ist merkwürdig, dass unter dieser Voraussetzung, die freilich nicht der Fall der Natur ist, die Schwere an der Oberfläche der Erde, sei es auf dem Meere, oder auf einer Insel oder dem festen Lande, von der Gestalt des über das Meer hervorragenden festen Landes unabhängig wird. --- Nach einer von dem Verf. später gemachten Bemerkung, ist dies Resultat von der Voraussetzung, dass das Erdsphäroid \emph{dieselbe} Dichtigkeit, wie das Meer, habe, unabhängig und setzt \emph{blos} die gleichförmige Dichtigkeit des erstern voraus. Und da die für die Schwere gefundene Formel eine merklich geringere Zunahme, vom Aequator nach dem Pol, gibt, als die Erscheinungen über die Pendelschwingungen, so wird dadurch bewiesen, dass die Dichtigkeit der Erde, \SourceLeaf{621}von der Oberfläche nach dem Mittelpunkt, zunimmt; es ist selbst nothwendig, dass diese Zunahme der Dichtigkeit sich bis zu einer bedeutenden Tiefe erstrecke. --- Ueber die Stellungen der Fixsterne, von welchen in den Ephemeriden die Mondsdistanzen angesetzt werden, von \textsc{Bouvard}. Aus dem \textsc{Bessel}-\textsc{Bradley}'schen Catalog für 1755 und dem \textsc{Piazzi}'schen für 1800 (indem jedoch die Rectascensionen des letztern, wie angegeben wird, nach Pariser Beobachtungen, ein wenig abgeändert sind) werden die Stellungen auf das Jahr 1830 übertragen, diese Reduction scheint jedoch nicht durchgängig genau gerechnet zu sein. --- Untersuchungen über vier ältere Cometen, von \textsc{Burckhardt}. Für den zweiten Cometen von 1766 findet Hr. \textsc{Burckhardt} mit Zuziehung der auf der Insel Bourbon von \textsc{Lanux} gemachten Beobachtungen eine elliptische Bahn mit einer Umlaufszeit von 5 Jahren; die Bewegungen der Cometen von 1774, 1723, 1729 lassen sich besser durch hyperbolische, als durch parabolische Bahnen darstellen. --- Von demselben Astronomen eine Tafel zur Bestimmung der Zeit, bei der parabolischen Bewegung, aus zwei Abständen von der Sonne und dem dazwischen enthaltenen Winkel; ferner kritische Bemerkungen über verschiedene von \textsc{Bradley} nur einmal beobachtete Sterne. --- Ueber die Englische Expedition nach dem Nordpol. --- Auszug aus der Beschreibung der Expedition des Capitain \textsc{Tuckey} zur Untersuchung des Flusses Zaire. --- Ueber die Gestalt der Erde, und das Gesetz der Schwere an ihrer Oberfläche von \textsc{Laplace}. Das hier abgedruckte enthält eine Uebersicht der Hauptresultate einer Abhandlung, welche der Verfasser in der Königl. Academie vorgelesen hat, und wovon ein vollständiger Abdruck vor uns liegt. Die Aufgabe ist darin auf eine vollständigere und naturgemässere Art behandelt, als bisher je geschehen war, und sowohl die Neuheit und Fruchtbarkeit der darin gebrauchten Methoden, als die Wichtigkeit der Resultate für die Geologie, machen sie zu einer der schönsten Arbeiten dieses grossen Geometers. --- Auszug eines Briefes des Hrn. W. \textsc{Lambton} an Hrn. \textsc{Delambre}, betreffend die Ostindische Gradmessung, die damals (im September 1817) bereits einen Meridianbogen von fast 10 Grad umfasste und noch immer weiter fortgesetzt wurde. Hr. \textsc{Lambton} gibt die Breite des südlichen Endpunkts Runnä zu $8^\circ\,9'\,38''4$, die des nördlichen Darumgedda zu $18^\circ\,3'\,23''26$, und die Länge des Bogens zu 598610 Fathoms bei $62^\circ$ Fahrenheit an und findet durch Vergleichung mit der Französischen Gradmessung die Abplattung $\frac{1}{309.5}$, aus der Englischen $\frac{1}{313.54}$, aus der Schwedischen $\frac{1}{307.15}$; hingegen aus Vergleichung des ganzen Bogens mit dem ganzen Bogen von Dünkirchen bis Montjouy, nach einer Formel von \textsc{Playfair} $\frac{1}{355}$ (diese starke Verschiedenheit hat ihren Grund in einer unrichtigen Verwandlung des Französischen Bogens aus Toisen in Fathoms). --- Auszug einiger Ortsbestimmungen in Schweden aus den Abhandlungen der Stockholmer Academie. --- Ueber den Cometen von 1818 von \textsc{Nicollet}. --- Auszug aus \textsc{Lemonniers} Beobachtungsjournal, in dem sich nicht weniger als 12 Uranusbeobachtungen finden, die erste vom 14. October 1750, die letzte am 18. December 1771. Im Jahr 1769 hatte dieser Astronom in 8 Tagen, vom 15. bis 23. Januar, den Planeten 6 mal unbewusst beobachtet, und man begreift kaum, wie ihm selbst die Entdeckung entgangen sein kann. Hr. \textsc{Bouvard}, welcher uns mit diesen Beobachtungen bekannt macht, ist beschäftigt, neue Uranus-Tafeln zu construiren. Uebrigens haben die Beobachtungen selbst bei der Beschaffenheit des Instruments, der Uhr und der geringen Sorgfalt des Beobachters keinen grossen Werth. --- Ueber die Refraction bei geringen Höhen von \textsc{Groombridge}. Dies ist ein Nachtrag zu einer frühern Abhandlung dieses Astronomen, in welcher er bei jedem einzelnen Stern nur das Mittel der Abweichung von seinen Tafeln angesetzt hatte; hier werden auch die äussersten Unterschiede mitgetheilt. Dies kann zwar einigermassen zur Beurtheilung des Grades der Uebereinstimmung der Beobachtungen unter sich dienen; allein in Beziehung auf die wichtige Frage, über welche die Meinungen der Astronomen bisher getheilt sind, ob es angemessen sei, bei der Reduction das äussere oder das innere Thermometer oder ein Mittel aus beiden anzuwenden, wäre zu wünschen, dass alle einzelnen Beobachtungen bekannt gemacht würden. Noch sind die Beobachtungen von vier Solstitien, von 1816---\SourceLeaf{622}1818 beigefügt; das Mittel der zwei Bestimmungen der Schiefe der Ekliptik auf den zwei Sonnensolstitien weicht hier nicht merklich von dem Mittel aus den zwei Wintersolstitien ab, und das Mittel aus allen gibt, für die Zeit der Herbstnachtgleiche von 1817, die mittlere Schiefe $23^\circ\,27'\,48''29$.
